\section{Related Work}

We review three lines of work that are most relevant to Subspace-SCAFFOLD: (i) classical federated learning algorithms and variance-reduction techniques, (ii) communication compression with error feedback, and (iii) subspace and low-rank methods for communication- and memory-efficient federated learning. We emphasize where existing methods fall short in simultaneously addressing heterogeneity, communication, and memory, which motivates our subspace-compatible SCAFFOLD design.

\subsection{Federated Learning and Variance Reduction}

The standard federated learning (FL) protocol is based on model averaging. McMahan et al.~\cite{mcmahan2017communication,mcmahan2017fedavg} introduced the FedAvg algorithm, which performs several local stochastic gradient steps on each participating client and periodically averages local models on the server. FedAvg is communication-efficient relative to centralized SGD, but under non-IID data it suffers from ``client drift'': local models overfit to their own data and the aggregated update can deviate significantly from the global descent direction.

To mitigate the effect of heterogeneity, a large body of work has focused on modifying the local objective or using variance reduction. FedProx~\cite{li2018federated} adds a proximal regularizer to constrain local updates around the current global model, improving convergence under system and data heterogeneity. Subsequent work has refined the theoretical foundations and practical aspects of heterogeneous FL~\cite{li2020federated}, but these methods still operate in the full parameter space and maintain full-dimensional model states.

SCAFFOLD~\cite{karimireddy2020scaffold} tackles client drift more directly via control variates. Each client and the server maintain control variate vectors that estimate the difference between the local and global gradients. By correcting local updates with these control variates, SCAFFOLD achieves significantly reduced variance and faster convergence under non-IID data, with strong theoretical guarantees. A growing literature explores variance-reduced FL beyond SCAFFOLD, including methods based on model discrepancy~\cite{zhang2025modeldiscrepancy}, proportional variance reduction under partial participation~\cite{wang2025proportionalvr}, and quantized variance-reduction schemes~\cite{wang2025quantizedvr}. Rodio~\cite{rodio2025variance} provides a recent overview of the many forms of variance reduction in FL.

However, a fundamental limitation shared by FedProx, SCAFFOLD, and most variance-reduced FL algorithms is that they rely on full-dimensional model updates and full-dimensional auxiliary states (e.g., control variates). For modern over-parameterized models, this implies that each client must store both the model parameters and additional $\mathcal{O}(d)$-dimensional statistics, which is prohibitive under tight memory budgets. Moreover, these methods are largely orthogonal to recent advances in subspace and low-rank training; naively combining them can break their theoretical assumptions, especially when the update subspace changes over time.

\subsection{Communication Compression and Error Feedback}

Orthogonal to algorithmic variance reduction, a rich line of work attacks the communication bottleneck in distributed and federated optimization via gradient compression. Classical approaches include sparsification, quantization, and their combinations. Sparsified SGD with memory~\cite{stich2018sparsesgd} and related schemes such as ATOMO~\cite{wang2018atomo} retain only a subset of coordinates or atomic components per iteration. Qsparse-Local-SGD~\cite{basu2020qsparse} combines aggressive sparsification, quantization, and local computation to reduce the communication volume while preserving linear-speedup guarantees under suitable assumptions.

Purely biased compression can severely harm convergence, particularly under non-IID data. Error feedback (EF) has emerged as a powerful generic mechanism to restore convergence guarantees by maintaining and re-injecting the compression error. Karimireddy et al.~\cite{karimireddy2019error} showed that EF can provably fix signSGD and many other biased compression schemes. Their analysis has inspired many EF-based algorithms in both centralized and federated settings, including contractive error feedback~\cite{li2024contractiveef} and accelerated EF-based distributed optimization~\cite{gao2025accelerated}. Recent work has further refined EF to better handle data heterogeneity and communication constraints: Jin et al.~\cite{jin2024magnitude} propose magnitude-aware sparsification with EF to fix signSGD under heterogeneity; Bao et al.~\cite{bao2025efskip} develop EFSkip, which carefully skips certain EF updates while retaining linear speedup; and Medjadji et al.~\cite{medjadji2025fedsparq} introduce FedSparQ, an adaptive sparse-quantization scheme with EF tailored to FL.

In the FL context, a number of works study gradient compression and EF on realistic wireless and edge systems. Zhang et al.~\cite{zhang2025gradient} couple gradient compression with correlation modeling for wireless traffic prediction, while Huang et al.~\cite{huang2025energy} design an energy-efficient FL method that jointly optimizes compression and training parameters. Quantized variance-reduction methods~\cite{wang2025quantizedvr} explicitly co-design the compressor and variance-reduced update rules for communication-limited wireless edge networks.

Despite their strong communication benefits, EF-based methods introduce new memory costs: each client must maintain a persistent error accumulator in $\mathbb{R}^d$. For very large models, the EF buffer can dominate the memory footprint, especially on low-resource devices. Moreover, EF is typically formulated in the full parameter space; its convergence guarantees depend on repeatedly feeding back the accumulated full-dimensional error. When updates are constrained to a changing low-dimensional subspace, it is unclear how to define or maintain an error memory that remains consistent across subspace switches. This tension between EF-style compression and subspace methods is not addressed by current algorithms.

\subsection{Subspace and Low-Rank Methods in Federated Learning}

Recent progress on parameter-efficient and memory-efficient training has highlighted the benefit of constraining updates to low-dimensional subspaces. In the centralized setting, low-rank gradient projection methods such as GaLore~\cite{zhao2024galore} project gradients onto a learned low-rank subspace, drastically reducing optimizer-state memory while maintaining competitive performance on large language models. Surveys on parameter-efficient fine-tuning (PEFT)~\cite{han2024peftsurvey} and federated low-rank adaptation~\cite{yang2025fedlora_survey} document a broad design space of low-rank adapters, sparsity structures, and optimization strategies.

In federated learning, low-rank and subspace ideas have been applied to reduce communication, computation, and memory simultaneously. Zhang et al.~\cite{zhang2025fedsub} propose FedSub, an efficient subspace algorithm for FL on heterogeneous data that restricts local updates to a carefully constructed low-dimensional subspace, alleviating client drift while greatly reducing per-round communication and local computation. Wu et al.~\cite{wu2025dpfedsub} extend this line with DPFedSub, which incorporates randomized subspace descent into a differentially private FL framework.

Another major thread is federated low-rank adaptation (FedLoRA) for fine-tuning large foundation models. Wang et al.~\cite{wang2024flora} and Koo et al.~\cite{koo2024robustfedlora} study FedLoRA variants that handle heterogeneous client capacities and data distributions while keeping communication and local memory low. Bian et al.~\cite{bian2024lorafair} refine aggregation and initialization of LoRA adapters, and Sun et al.~\cite{sun2024improvinglora} focus on improving LoRA in privacy-preserving FL. Zhou et al.~\cite{zhou2025ilora} introduce ILoRA, which explicitly addresses instability from misaligned client subspaces and rank incompatibility during aggregation. Peng et al.~\cite{peng2025fedhl} propose FedHL to achieve unbiased aggregation for heterogeneous low-rank adapters. More broadly, Yang et al.~\cite{yang2025fedlora_survey} and Han et al.~\cite{han2024peftsurvey} summarize how low-rank and subspace techniques can drastically reduce the memory and communication footprints of distributed training.

These works demonstrate that subspace and low-rank parameterizations are highly effective at reducing both communication and memory in FL. Yet, they are largely developed independently of SCAFFOLD-style variance reduction and EF-based compression. A key technical challenge is that control variates and EF memories are defined in the full parameter space and rely on temporal consistency of the update direction. When clients switch subspaces over time---as in adaptive low-rank projection methods---historical variance-reduction statistics may become misaligned, leading to biased corrections or even divergence.

To the best of our knowledge, existing subspace FL methods such as FedSub~\cite{zhang2025fedsub} and DPFedSub~\cite{wu2025dpfedsub}, as well as FedLoRA-style approaches~\cite{wang2024flora,zhou2025ilora,koo2024robustfedlora}, do not provide a principled mechanism to retain the benefits of SCAFFOLD-like control variates or EF under dynamically evolving subspaces. Conversely, existing variance-reduced and EF-based FL algorithms~\cite{karimireddy2020scaffold,karimireddy2019error,wang2025quantizedvr,zhang2025modeldiscrepancy,wang2025proportionalvr} do not account for subspace-constrained updates and assume full-dimensional auxiliary states. This mismatch leaves a gap: we lack methods that are simultaneously (i) communication-efficient, (ii) memory-efficient, and (iii) robust to data heterogeneity via variance reduction when updates are restricted to low-dimensional, possibly time-varying subspaces. Subspace-SCAFFOLD is designed precisely to fill this gap by redefining control variates in a subspace-aware manner and ensuring their compatibility with subspace switches.

